% !PRINTED 453
\weberpage{453}

%<<CONT
the $P_8$ may fittingly be called the \gesperrt{principal poles}, and the axes passing through them may be denoted by $A_8, A_6, A_4, A_2$.

To be singled out further are the 28 pairs of connecting lines of a $P_6$ with the two associated $P_3$, which we may denote by $A_3$, and finally the 24 connecting lines of each two $P_7$ belonging together, as whose representatives the sides of the coordinate triangle are to be regarded, and which may be denoted by $A_7$.

\begin{center}
\S.\,121.\\
\gesperrt{Invariant curves of the group $G_{168}$.}
\end{center}

A form of degree $\mu$, $\varphi\,(x_1, x_2, x_3)$, will in general pass over, under the substitutions of the group, into 168 different forms. For special forms of $\varphi$, however, this number can diminish, and then the substitutions under which $\varphi$ remains unchanged form a group $G'$ contained in $G_{168}$, whose degree must be a divisor of 168. The number of forms into which $\varphi$ passes over is the index of the divisor $G'$, and each of these different forms of $\varphi$ remains unchanged under a group conjugate to $G'$. The form $\varphi$ is an absolute invariant of the group $G'$.

The equation $\varphi = 0$ represents a curve referred to the coordinate triangle $x_1, x_2, x_3$, and this curve is mapped by the substitutions of $G_{168}$ onto other curves. If the 168 image curves are not all different from one another, then the curve $\varphi$ remains unchanged under the substitutions of a group $G'$, or changes only by a constant factor, and is thus a relative invariant of $G'$. A straight line remains unchanged under substitutions other than the identical one only when it belongs to the axes described in the preceding section.

All properties and relations between points, lines and curves that are indestructible under linear transformation (the so-called projective properties of geometry) are preserved in the images. If therefore, e.\,g., a straight line is a tangent or inflectional tangent or double tangent of a curve, then all images of the straight line stand in the same relation to the images of the curve. Likewise, if a point
%CONT>>
